% vim: ai sw=2 spell spelllang=cs

\begin{frame}
  \frametitle{Profilování}

  \heading{Co zabírá nejvíce času při vykonávání?}
  \begin{enumerate}
    \item \cmd{gprof}
    \begin{itemize}
      \item Při překladu: \cmd{gcc -pg ...}
      \item Potom se program spustí a data se uloží na disk
      \item \cmd{gprof} data načte a zobrazí
%      \item Demo: dump
    \end{itemize}

    \pause

    \item \cmd{perf}
    \begin{itemize}
      \item Používá čítače v procesoru (nutná podpora CPU)
      \item Zahrnuje ve výpisech chování jádra
      \item Umí pracovat s \cmd{gcc -pg} binárkami
      \item Spuštění: \cmd{perf record -g -- binarka --volby}
      \item Dokumentace: \url{http://perf.wiki.kernel.org}
    \end{itemize}

  \end{enumerate}

\end{frame}

\begin{frame}
  \frametitle{Úkol}

  \heading{Profilování}
  \begin{enumerate}
    \item Projděte si \iffimuni \file{pb173-bin/10/prof.c} \else XXX \fi
    \item Spusťte \file{prof_pg}
    \item Prostudujte výstup \cmd{gprof -b prof_pg}

    \bigskip

    \item Nahrajte si výstup \cmd{perf}
    \begin{itemize}
      \item \cmd{perf record -g prof}
    \end{itemize}
    \item Zobrazte si výstup
    \begin{itemize}
      \item \cmd{perf report -g}
    \end{itemize}
    \item Zopakujte \cmd{perf} pro \file{prof_pg}
  \end{enumerate}

\end{frame}

